\section{Questions ouvertes et expériences discriminantes}
\label{sec:perspectives}
Le premier résultat à prolonger est la réponse collective à la technologie ;
le partage du surplus constitue une seconde voie de recherche. Les pistes
suivantes sont des intentions scientifiques, pas des résultats acquis ni
des protocoles définitivement arrêtés.

\paragraph{Hétérogénéité du trio $(K_0,A,\gamma)$.}
Il faut étudier ensemble la dotation initiale, la productivité et la concavité :
trio identique à toutes les naissances, tirage dans une distribution bornée,
corrélations entre composantes, ou dépendance à l'état de la simulation.
Comparer les conventions à $\kappa_0$ contrôlé puis libre permettrait de séparer
hétérogénéité technologique et déplacement d'échelle. Si $\gamma_i$ varie,
la dimension de $A_i$ varie aussi : un tirage uniforme de valeurs numériques
de $A_i$ n'a pas de sens indépendant de l'unité choisie. Il faut fixer une
échelle commune $K_{\rm ref}$ et, par exemple, spécifier
$a_i=A_i K_{\rm ref}^{\gamma_i-1}$, puis documenter la loi jointe de
$(K_{0,i}/K_{\rm ref},a_i,\gamma_i)$.
Les sorties à suivre incluent survie des cohortes, production par entité,
population, réseau de crédit et distributions de bilan ; quels indicateurs
endogènes résument leur évolution micro- et macroscopique reste une question.

\paragraph{Causalité entre pas et hypothèse de criticité.}
Les avalanches enregistrées ferment la cascade \emph{dans un pas}.
Cette convention de mesure ne signifie pas que la causalité du système
s'y limite : le carnet de dettes, les stocks et les cohortes persistent et
modifient les risques futurs. Une entité classée « racine » aujourd'hui peut
donc porter une fragilité accumulée hier. Inversement, la criticité n'est pas
définie en général comme un phénomène limité à un seul pas.
La SOC reste ouverte : il faudrait définir une réponse spatio-temporelle à
une perturbation, des horizons, un contrefactuel et une analyse de taille finie.
Ni $b_1$ intra-pas, ni une corrélation retardée, ni une droite log--log ne
suffisent à trancher seuls.

\paragraph{Sur-altruisme et temporalité du contrat.}
L'extension proposée consiste à accepter $p<0$ dans $rq=L+p\Delta$ :
la prêteuse reçoit initialement moins que sa perte productive de référence.
La dissipation et l'évolution des deux entités pourraient changer ce bilan
ultérieurement ; cette inversion est une hypothèse à tester, pas une
conséquence automatique du signe de $p$.\footnote{Le moteur de la campagne
impose actuellement $0\leq p\leq1$ pour la règle \texttt{bargain}.
Un prolongement devra spécifier l'admissibilité du service $rq$, le maintien
ou non du taux nominal, l'horizon, les faillites et le contrefactuel sans prêt.
Un gain productif instantané négatif n'est pas une valeur actualisée nette
négative : le modèle ne fixe pas de taux d'actualisation.}
Un premier test isolerait une paire, avec et sans contrat, puis répéterait
la comparaison dans une population ; seuls les seconds essais renseigneraient
les rétroactions démographiques et la propagation.

\paragraph{Distributions et comparaison empirique.}
Pareto demeure une hypothèse suggérée par les graphiques et imparfaitement
testée faute de temps. Le corps exponentiel n'est plus une cible obligatoire :
l'objectif est d'identifier la forme et le mécanisme, non de reproduire un
gabarit universel supposé. Comparer les familles sur le même échantillon et le
même domaine exige de tenir compte de leur flexibilité.\footnote{Pour deux
familles emboîtées ajustées correctement, la vraisemblance maximale ne diminue
pas lorsqu'on ajoute un paramètre ; pour des familles non emboîtées, le seul
nombre de paramètres ne classe pas les ajustements. Une comparaison par
AIC/BIC, validation hors échantillon ou test adapté complète l'adéquation
absolue. Les observations de plusieurs instantanés d'un même run ne sont pas
des réplications indépendantes.}
La priorité empirique est une distribution de bilan ou d'intérêts dont
l'unité statistique, la fenêtre et les conventions comptables soient
compatibles ; l'entité générique ne sera assimilée à un acteur réel qu'à
cette condition. Robustesse des queues, seuils de coupure, limite continue
et autres règles de crédit restent des expériences possibles.
